\documentclass[11pt,a4paper]{article}
\usepackage[utf8]{inputenc}
\usepackage[T1]{fontenc}
\usepackage[brazil]{babel}
\usepackage[a4paper,margin=2.3cm]{geometry}
\usepackage{enumitem}
\usepackage{hyperref}
\usepackage{xcolor}
\setlength{\parindent}{0pt}
\setlength{\parskip}{5pt}
\sloppy
\setlength{\emergencystretch}{3em}
\hypersetup{colorlinks=true,urlcolor=blue,linkcolor=black}
\newcommand{\cabecalho}[2]{\begin{center}{\small \textbf{Sprint 2 -- #1}}\\[1mm]\textit{Inventário de Ativos e Vulnerabilidades}\\Matheus Moreli Durão Silva -- 12621CBS217\\Cibersegurança -- Bacharelado -- Noturno\end{center}\hrule\vspace{4mm}}
\newcommand{\bloco}[1]{\par\vspace{2mm}\textbf{#1}\par}
\newcommand{\codigo}[1]{\texttt{#1}}

\begin{document}\cabecalho{S2\_13 -- Cadastro de Vulnerabilidades}{}
\bloco{Dados do cadastro} O registro possui CVE, CWE, CVSS, descrição, fonte, data da fonte, impacto no ativo, prioridade, tratamento, status e verificação. Ele é guardado dentro do ativo afetado, evitando uma vulnerabilidade sem relação com o inventário.
\bloco{Validações} \codigo{validar\_cve} separa o texto por hífen e aceita apenas \codigo{CVE-AAAA-NÚMERO}. \codigo{validar\_cvss} usa \codigo{float} e exige 0.0 a 10.0. Antes de montar o objeto, \codigo{consultar\_por\_id} confirma que o ativo existe. Fonte e data são textos obrigatórios.
\bloco{Dois exemplos reais, ativos simulados} \codigo{CVE-2019-0708}, CWE-416: falha no Remote Desktop Services; cadastrar apenas se o ativo usa versão afetada do Windows e registrar patch/restrição de RDP como tratamento. \codigo{CVE-2021-44228}: falha de JNDI no Log4j2; cadastrar apenas após confirmar \codigo{log4j-core} afetado, com atualização e verificação de dependências como tratamento. Fontes: \url{https://nvd.nist.gov/vuln/detail/CVE-2019-0708} e \url{https://nvd.nist.gov/vuln/detail/CVE-2021-44228}. O inventário demonstrativo não é prova de exposição real.
\bloco{Como testar} Os testes aceitam um CVE e CVSS válidos e rejeitam CVE malformado e CVSS 11. Após cadastro válido, a consulta exata pelo CVE encontra uma ocorrência.
\bloco{Checklist atendido} Formato CVE, limite CVSS, ativo existente, fonte verificável e cuidado com dados não confirmados.
\end{document}
